\chapter{P\&L Explain and Independent Price Verification}\label{ch:rc:pnl-explain-and-independent-price-verification}

Every evening a product controller must explain each dollar a desk made or lost: so much from
the move in rates, so much from volatility, so much from the passage of time, and a remainder
that nobody can account for. On most days the remainder is small. On the day a swaption book
made 4.57 million dollars after a 30 basis point jump in rates and a rise in volatility, the
desk's own explain accounted for 4.17 million, and 400 thousand dollars were unexplained. An
\omterm{def:rc:pnl-explain-and-independent-price-verification:unexplained}{unexplained P\&L} is either a missing risk, a wrong mark, or a booking error, and \omterm{def:rc:pnl-explain-and-independent-price-verification:pc}{product control}
must find out which. This chapter builds the tools: attribution by sensitivities and by
revaluation, the verification of marks against independent prices, the reserves and the
prudent-valuation adjustments that reflect what the marks cannot know, and the control
functions that run them. One Quant Book~1 introduced P\&L attribution for a trading firm; here it
is applied to a derivatives book, where the risks interact.

\section{Attribution: risk-based and revaluation-based}

\begin{definition}[Risk-based attribution]\label{def:rc:pnl-explain-and-independent-price-verification:risk}
\emph{Risk-based attribution}\index{risk-based attribution} explains the day's P\&L by a Taylor
expansion in the risk-factor moves, using the start-of-day Greeks:
$\Delta V\approx\delta\,\Delta f+\tfrac12\Gamma\,\Delta f^2+\nu\,\Delta\sigma+\tfrac12\,\mathrm{volga}\,\Delta\sigma^2+\mathrm{vanna}\,\Delta f\,\Delta\sigma+\theta\,\Delta t$.
\end{definition}

\begin{definition}[Revaluation-based attribution]\label{def:rc:pnl-explain-and-independent-price-verification:reval}
\emph{Revaluation-based attribution}\index{revaluation-based attribution} moves the \omterm{def:rc:market-risk-measures:factor}{risk factors}
from their start-of-day to their end-of-day values one at a time, revaluing in full after each, and
assigns each step's change to its factor. The steps add up to the P\&L exactly; their split depends
on the order.
\end{definition}

\begin{definition}[Unexplained P\&L]\label{def:rc:pnl-explain-and-independent-price-verification:unexplained}
The \emph{unexplained P\&L}\index{unexplained P\&L} is the difference between the actual P\&L and
the sum of the attributed parts. Controllers set thresholds on it, in money and as a share of the
P\&L, beyond which it must be investigated.
\end{definition}

\begin{example}[A risk reversal on a quiet day]\label{ex:rc:pnl-explain-and-independent-price-verification:book}
An illustrative book is long USD~320 million of one-year-into-ten-year payer swaptions struck at 4.50\%
and short the same notional of receivers at 3.50\%, a risk reversal around a forward of 4.00\%, marked at
a normal volatility of 95 basis points (0.75 years to expiry, annuity 8.0). The two legs are worth the
same, so the book is worth zero, and by symmetry its gamma, vega and theta are zero too: it is a pure
delta position, 139\,103 dollars per basis point. Its cross sensitivity is not zero: the vanna is 1\,086
dollars per basis point of rate and basis point of volatility, because a higher rate moves the long payer
towards the money and the short receiver away from it.
\end{example}

\omcode{code/firm/pnlexplain/firm_pnlexplain.py}{19}{41}{Finite-difference Greeks including the cross term, and the risk-based attribution with or without it.}\label{lst:rc:pnl-explain-and-independent-price-verification:risk}

\section{Unexplained P\&L and what it reveals}

\begin{example}[The big-move day]\label{ex:rc:pnl-explain-and-independent-price-verification:day}
The forward rises 30 basis points and the volatility 12. The book makes USD~4\,571\,501. The desk's
explain, delta, gamma, vega and theta, gives 4\,173\,095, all from delta: 398\,406 dollars are unexplained,
9\% of the day. Adding the vanna term, $1\,086\times30\times12$, explains 391\,087 of them; 7\,319 remain,
higher-order terms (\cref{fig:rc:pnl-explain-and-independent-price-verification:attribution}). The
\omterm{def:rc:pnl-explain-and-independent-price-verification:unexplained}{unexplained P\&L} grows with the product of the rate and volatility moves, which is why it appears on big days
(\cref{fig:rc:pnl-explain-and-independent-price-verification:unexplained}).
\end{example}

\begin{omfigure}
\begin{tikzpicture}
\begin{axis}[width=11.5cm,height=5.2cm,ybar=4pt,bar width=24pt,
  symbolic x coords={delta,vanna,unexplained},xtick=data,enlarge x limits=0.25,
  ylabel={USD thousand},
  tick label style={font=\small},label style={font=\small},
  ymin=0,ymax=4800,grid=major,grid style={gray!25},
  nodes near coords,every node near coord/.append style={font=\footnotesize,/pgf/number format/fixed,/pgf/number format/precision=0},
  legend columns=2,legend style={font=\footnotesize,at={(0.5,-0.2)},anchor=north,draw=none,fill=none,/tikz/every even column/.append style={column sep=8pt}},
  table/col sep=comma]
\addplot[fill=omExo!40,draw=omExo] table[x=term,y=desk]{figdata/rates-credit-risk/27-pnl-explain-and-independent-price-verification/attribution.csv};
\addplot[fill=omDef!60,draw=omDef] table[x=term,y=full]{figdata/rates-credit-risk/27-pnl-explain-and-independent-price-verification/attribution.csv};
\legend{desk's explain,with the cross term}
\end{axis}
\end{tikzpicture}

\omcaption{Attribution of the risk reversal's P\&L on the big-move day (USD~4.57 million): the desk's
explain leaves 398 thousand dollars unexplained; the vanna term accounts for almost all of it. Data: the
chapter's tutorial.}\label{fig:rc:pnl-explain-and-independent-price-verification:attribution}
\end{omfigure}

\begin{omfigure}
\begin{tikzpicture}
\begin{axis}[width=11.5cm,height=5.2cm,
  xlabel={move of the forward rate (bp), volatility up 12 bp},ylabel={unexplained (USD thousand)},
  tick label style={font=\small},label style={font=\small},
  xmin=-40,xmax=40,ymin=-650,ymax=650,grid=major,grid style={gray!25},
  legend columns=2,legend style={font=\footnotesize,at={(0.5,-0.25)},anchor=north,draw=none,fill=none,/tikz/every even column/.append style={column sep=8pt}},
  table/col sep=comma]
\addplot[omThm,very thick,mark=*,mark size=1.2pt] table[x=bp,y=desk]{figdata/rates-credit-risk/27-pnl-explain-and-independent-price-verification/unexplained.csv};
\addplot[omDef,very thick,mark=square*,mark size=1.2pt] table[x=bp,y=full]{figdata/rates-credit-risk/27-pnl-explain-and-independent-price-verification/unexplained.csv};
\legend{desk's explain,with the cross term}
\end{axis}
\end{tikzpicture}

\omcaption{\omterm{def:rc:pnl-explain-and-independent-price-verification:unexplained}{Unexplained P\&L} of the risk reversal as a function of the day's rate move, with the
volatility up 12 basis points. Without the cross term the unexplained part is proportional to the rate
move; with it, only third-order terms remain. Data: the chapter's tutorial.}\label{fig:rc:pnl-explain-and-independent-price-verification:unexplained}
\end{omfigure}

\begin{example}[The same day by revaluation]\label{ex:rc:pnl-explain-and-independent-price-verification:reval}
Moving the forward first, then the volatility, then time: the forward contributes USD~4\,240\,447, the
volatility 336\,046, time $-4\,992$. Moving the volatility first: the volatility contributes nothing (the
book's vega is zero at the start), the forward 4\,576\,493. Both add up to the P\&L; the cross effect
lands in whichever factor moves second.
\end{example}

\begin{definition}[Product control]\label{def:rc:pnl-explain-and-independent-price-verification:pc}
\emph{Product control}\index{product control} is the finance function, independent of the desks, that
produces and explains the daily P\&L, verifies the marks, sets reserves and reconciles the \omterm{def:rc:regulatory-capital-for-trading-books:books}{trading
books} with the general ledger.
\end{definition}

\section{Independent price verification}

\begin{definition}[Independent price verification]\label{def:rc:pnl-explain-and-independent-price-verification:ipv}
\emph{Independent price verification}\index{independent price verification} (IPV) is the process by
which the market prices and model inputs used by a desk are verified against sources independent of
it (broker quotes, trades, consensus services) and the marks corrected when they fall outside a
tolerance.
\end{definition}

In a consensus service, dealers submit their mid marks for a list of contracts once a month and receive
the average of the accepted submissions; a dealer whose submission is rejected does not receive the
consensus. The leading service, Totem, began in 1997 with six dealers and had about 120 participants and some 30
submitters per contract, according to a 2020 study of its data.

\omcode{code/firm/pnlexplain/firm_pnlexplain.py}{56}{62}{IPV: a mark outside the consensus band is moved to the nearest boundary.}\label{lst:rc:pnl-explain-and-independent-price-verification:ipv}

\begin{example}[Verifying the volatility marks]\label{ex:rc:pnl-explain-and-independent-price-verification:ipv}
The desk marks both strikes at 95 basis points; an illustrative consensus gives 97 for the 4.50\% payer
and 91 for the 3.50\% receiver, the market's skew, with a tolerance of 2 basis points. The payer's mark
sits on the boundary and stands; the receiver's is moved to 93, which lowers the value of the short
receivers and adds USD~146\,482. The flat marks also explain the vanna: a desk that marks no skew sees
no reason for its vega to change when rates move.
\end{example}

\begin{dated}{2026-09}{Marks at the boundary}\label{dat:rc:pnl-explain-and-independent-price-verification:cio}
JPMorgan's 2013 task force report describes its Chief Investment Office's price testing: each trader
mark was compared with a mid price computed by the valuation control group and a threshold reflecting
the bid--offer spread, and marks outside were to be moved to the nearest boundary. In the March 2012
testing many positions were marked at or near the boundary of the bid--offer spread; being within the
thresholds, they were accepted.
\end{dated}

\section{Reserves and the fair-value hierarchy}

\begin{definition}[Fair-value hierarchy]\label{def:rc:pnl-explain-and-independent-price-verification:hierarchy}
The \emph{fair-value hierarchy}\index{fair-value hierarchy} classifies fair-valued positions by the
observability of their inputs (IFRS~13): level~1, unadjusted quoted prices in active markets for
identical instruments; level~2, other directly or indirectly observable inputs; level~3, unobservable
inputs. Level~3 positions need the
most verification and carry the largest reserves.
\end{definition}

Reserves adjust fair values for what a mid mark ignores: the bid--offer cost of exiting, the
uncertainty of the model and its parameters, and day-one profits that cannot yet be recognised
(One Quant Book~5, chapter~27). They are part of fair value; regulators ask for more.

\section{Prudent valuation}

\begin{definition}[Prudent valuation, additional valuation adjustment]\label{def:rc:pnl-explain-and-independent-price-verification:pv}
\emph{Prudent valuation}\index{prudent valuation} values fair-valued positions at the price at which
the bank is 90\% confident it could exit them. The \emph{additional valuation
adjustments}\index{additional valuation adjustment} (AVAs) are the differences between fair value and
prudent value, by category (market price uncertainty, close-out costs, \omterm{def:rc:model-risk-and-validation:risk}{model risk} and others); they are
deducted from common equity.
\end{definition}

\begin{example}[The AVA of the book]\label{ex:rc:pnl-explain-and-independent-price-verification:ava}
With contributors' volatilities dispersed by 2.0 and 2.5 basis points around the consensus, the market
price uncertainty AVA, fair value less the 90\% confidence exit value, is USD~190\,578 for the payers and
232\,964 for the receivers; summed after the 50\% aggregation factor, 211\,771. The simplified approach,
0.1\% of the absolute fair values (USD~6.87 million), would give 6\,872: for an option book whose value is
small against its risk, the simplified figure understates the uncertainty.
\end{example}

\begin{dated}{2026-09}{Prudent valuation in the EU}\label{dat:rc:pnl-explain-and-independent-price-verification:pv}
Under the EU's \omterm{def:rc:pnl-explain-and-independent-price-verification:pv}{prudent valuation} standard (Delegated Regulation 2016/101), the prudent value is the value
at which a bank is 90\% confident of exiting; institutions with less than EUR~15 billion of absolute fair-valued
assets and liabilities may use the simplified approach (0.1\% of that sum), others the core approach, in
which the market price uncertainty, close-out cost and \omterm{def:rc:model-risk-and-validation:risk}{model risk} AVAs are summed after a 50\% aggregation
factor (temporarily 66\% after the 2020 amendment for extreme volatility). Total AVAs are deducted from
common equity tier~1.
\end{dated}

\section{Tutorial: explaining a day}\label{tut:rc:pnl-explain-and-independent-price-verification:tutorial}

\begin{tutorial}
\textbf{Goal.} Explain the risk reversal's big-move day by Greeks and by revaluation, verify its marks and
compute its prudent-valuation adjustment. \textbf{End state:} the numbers of
\cref{ex:rc:pnl-explain-and-independent-price-verification:book,ex:rc:pnl-explain-and-independent-price-verification:day,ex:rc:pnl-explain-and-independent-price-verification:reval,ex:rc:pnl-explain-and-independent-price-verification:ipv,ex:rc:pnl-explain-and-independent-price-verification:ava}
and the charts.
\begin{tutsteps}
\item \textbf{Greeks}: \texttt{greeks(value, START, BUMPS)}.
\item \textbf{Explain}: \texttt{day()} with and without the cross term; \texttt{revaluation()} in two orders.
\item \textbf{IPV and AVA}: \texttt{ipv\_table()}, \texttt{ava\_table()}.
\item \textbf{Charts}: \texttt{fig\_rc\_pnlexplain.py}.
\end{tutsteps}
\textbf{What to change next.} Mark the skew from the consensus and redo the day: the vanna becomes a vega
that the explain sees; add a second forward and a correlation cross term.
\end{tutorial}

\section{Build: the P\&L explain}\label{bld:rc:pnl-explain-and-independent-price-verification:pnlexplain}

\begin{build}
\bfield{Purpose} The firm's daily explain, price verification and \omterm{def:rc:pnl-explain-and-independent-price-verification:pv}{prudent valuation}, fed by the P\&L of the
position keeper of One Quant Book~1 (\texttt{firm.pnl}) and by the Greeks of the pricing library.
\bfield{Interface} \texttt{greeks(value, state, bumps)}; \texttt{risk\_based(g, df, ds, days, cross)};
\texttt{revaluation\_based(value, start, end, order)}; \texttt{ipv(marks, consensus, tolerance)};
\texttt{mpu\_ava}; \texttt{aggregate\_ava}; \texttt{simplified\_ava}.
\bfield{Rules} Greeks at the start of day; explain thresholds in money and share; IPV moves marks to the
tolerance boundary; 90\% confidence and 50\% aggregation for AVAs.
\bfield{Acceptance tests} \texttt{code/firm/pnlexplain/tests/}: a quadratic book is explained exactly with the
cross term and not without; revaluation steps add up and depend on the order; IPV boundaries; AVA of a normal
dispersion; aggregation and simplified figures.
\bfield{Stretch} Explain by trade and by \omterm{def:rc:market-risk-measures:factor}{risk factor} across a whole book; alerting on \omterm{def:rc:pnl-explain-and-independent-price-verification:unexplained}{unexplained P\&L}; close-out
cost and model-risk AVAs; consensus data ingestion.
\end{build}

\begin{omsources}
\item European Banking Authority, final draft RTS amending the RTS on prudent valuation, EBA/RTS/2020/04, 2020.
\item Basel Committee on Banking Supervision, \emph{Supervisory guidance for assessing banks' financial
instrument fair value practices}, April 2009.
\item JPMorgan Chase \& Co., \emph{Report of the Management Task Force Regarding 2012 CIO Losses}, 2013.
\item IFRS~13 \emph{Fair Value Measurement}, paragraphs 76, 81 and 86, as adopted in Commission Regulation (EU)
No 1255/2012.
\item L.~M.\ Ergun and A.~Uthemann, ``Higher-order uncertainty in financial markets: evidence from a
consensus pricing service'', Systemic Risk Centre Discussion Paper 98, June 2020.
\end{omsources}

\section{Exercises}

\begin{exercise}[$\star$]\label{exo:rc:pnl-explain-and-independent-price-verification:1}
Compute the vanna term of the big-move day from the book's vanna of 1\,086 dollars per basis point squared.
\end{exercise}

\begin{exercise}[$\star$]\label{exo:rc:pnl-explain-and-independent-price-verification:2}
Why are the risk reversal's gamma, vega and theta zero at the start of the day?
\end{exercise}

\begin{exercise}[$\star$]\label{exo:rc:pnl-explain-and-independent-price-verification:3}
A mark is 38 against a consensus mid of 35 with a tolerance of 2. What is the verified mark?
\end{exercise}

\begin{exercise}[$\star\star$]\label{exo:rc:pnl-explain-and-independent-price-verification:4}
Why does \omterm{def:rc:pnl-explain-and-independent-price-verification:reval}{revaluation-based attribution} depend on the order of the factors, and how would you remove the
dependence?
\end{exercise}

\begin{exercise}[$\star\star$]\label{exo:rc:pnl-explain-and-independent-price-verification:5}
Scale the \omterm{def:rc:pnl-explain-and-independent-price-verification:unexplained}{unexplained P\&L} to a day with a 15 basis point rate move and the same volatility move.
\end{exercise}

\begin{exercise}[$\star\star$]\label{exo:rc:pnl-explain-and-independent-price-verification:6}
Why is the market price uncertainty AVA of an option book large compared with its fair value?
\end{exercise}

\begin{exercise}[$\star\star\star$]\label{exo:rc:pnl-explain-and-independent-price-verification:7}
\emph{Coding.} Compute the AVA of the book with the contributors' dispersion doubled.
\end{exercise}

\begin{exercise}[$\star\star\star$]\label{exo:rc:pnl-explain-and-independent-price-verification:8}
\emph{Find the flaw.} ``All our marks are within the IPV thresholds, so our valuation is fine.''
\end{exercise}

\section{Problem: The Unexplained 400 Thousand}

\begin{problem}[{Weekend problem --- the missing cross term}]\label{pb:rc:pnl-explain-and-independent-price-verification:1}
The product controller of the risk reversal desk sees USD~398\,406 unexplained on the big-move day and must
report by morning.

\medskip\noindent\textbf{Part I --- The numbers.}
\begin{enumerate}
\item Give the actual P\&L and the desk's explain.
\item Give the unexplained part and its share of the P\&L.
\item Give the delta and vanna of the book.
\item Give the vanna term and the residual after it.
\item What is the residual made of?
\end{enumerate}

\medskip\noindent\textbf{Part II --- Diagnosis.}
\begin{enumerate}[resume]
\item Which three causes of \omterm{def:rc:pnl-explain-and-independent-price-verification:unexplained}{unexplained P\&L} would you check first, and in what order?
\item How does \omterm{def:rc:pnl-explain-and-independent-price-verification:reval}{revaluation-based attribution} locate the cross effect?
\item What does the volatility-first order show?
\item Why did the \omterm{def:rc:pnl-explain-and-independent-price-verification:unexplained}{unexplained P\&L} not appear on quiet days?
\item What does the desk's \omterm{def:rc:vanilla-rates-options:flat}{flat volatility} mark have to do with it?
\end{enumerate}

\medskip\noindent\textbf{Part III --- Controls.}
\begin{enumerate}[resume]
\item Give the IPV adjustment and its effect on the book's value.
\item Give the market price uncertainty AVA, as core and simplified approaches would compute it.
\item Where would you set the unexplained-P\&L threshold for this desk?
\item Should the desk be allowed to explain the day by revaluation only?
\item What would you report to the risk committee?
\end{enumerate}

\medskip\noindent\textbf{Part IV --- Judgement.}
\begin{enumerate}[resume]
\item When is an \omterm{def:rc:pnl-explain-and-independent-price-verification:unexplained}{unexplained P\&L} a sign of a booking error rather than a missing risk?
\item Why must \omterm{def:rc:pnl-explain-and-independent-price-verification:pc}{product control} be independent of the desk?
\item How does this connect with the FRTB attribution test of chapter~23?
\item State the \emph{named result}: the missing cross term and the attribution after it is added.
\item In one sentence: what does an \omterm{def:rc:pnl-explain-and-independent-price-verification:unexplained}{unexplained P\&L} tell you?
\end{enumerate}
\end{problem}

\section{Interview questions}

\begin{interviewq}[$\star$ \iqroles{risk, bank}]\label{iq:rc:pnl-explain-and-independent-price-verification:1}
Explain a day's P\&L on an option book. What goes into the explain?
\end{interviewq}

\begin{interviewq}[$\star\star$ \iqroles{trader, risk}]\label{iq:rc:pnl-explain-and-independent-price-verification:2}
What are the usual causes of \omterm{def:rc:pnl-explain-and-independent-price-verification:unexplained}{unexplained P\&L}?
\end{interviewq}

\begin{interviewq}[$\star\star$ \iqroles{bank}]\label{iq:rc:pnl-explain-and-independent-price-verification:3}
What is \omterm{def:rc:pnl-explain-and-independent-price-verification:ipv}{independent price verification}, and where does it get hard?
\end{interviewq}

\begin{interviewq}[$\star\star$ \iqroles{researcher, bank}]\label{iq:rc:pnl-explain-and-independent-price-verification:4}
What is \omterm{def:rc:pnl-explain-and-independent-price-verification:pv}{prudent valuation} and how does it differ from fair value?
\end{interviewq}

\begin{interviewq}[$\star\star\star$ \iqroles{developer}]\label{iq:rc:pnl-explain-and-independent-price-verification:5}
Design a daily P\&L explain system for a large derivatives book.
\end{interviewq}

\begin{interviewq}[$\star\star\star$ \iqroles{risk, bank}]\label{iq:rc:pnl-explain-and-independent-price-verification:6}
How would you detect marks that are systematically favourable to a desk?
\end{interviewq}
